\section{Conclusion}\label{sec:conclusion}

If local summaries preserve the evidence a task oracle needs, and local
merges preserve the cross-span interactions, then preference
optimization can be carried out on compression trees rather than full
documents. The neural-operator layer now makes explicit where the
learned architecture enters: discretization invariance and compact-set
approximation for the equation-(6) operator family explain how an ideal
operator can be realized to a chosen tolerance, the transformer-inclusion
argument identifies finite attention stacks with that architecture
surface, and the realization bridge of
Proposition~\ref{prop:neural-operator-bridge} converts that tolerance
into approximate local-law budgets. The oracle/projection objective
supplies the training interpretation: fit the task oracle while
weighting C1/C2/C3 as a projection toward the theorem-eligible
subspace. The projection/structured-error equivalence says this is the
same as assuming the approximation error is structured so that its zero
set is exactly local-law validity. Everything downstream then reuses
the existing theorem-backed route: the structural bundle gives root
preservation, the optimization bundle gives objective equivalence for
DPO/GRPO, and the certificate bundle turns sampled audits into the
finite-sample gap decomposition laid out in Section~\ref{sec:discussion}:
an empirical IPW estimate plus a fixed calibration residual plus an
estimation envelope that shrinks monotonically in the total number of
labels collected, global or local. The four statements the Discussion
uses (the main gap bound, fixed calibration, envelope rate, and
global/local decomposition) are cross-referenced to their formalized
versions in Appendix~\ref{app:proof-artifacts}.

The argument runs from an algebraic target to a measured application.
Section~\ref{sec:mergeable-sketches} grounds the framework in
the data-systems literature on mergeable sketches, with HLL as the
empirical representative and a broader encode/merge/decode class
covering the rest; Section~\ref{sec:framework} states the local laws
that lift that discipline to a learned summarizer, and
Section~\ref{sec:theorems} the preservation, equivalence, and
certificate results built on them. Two controlled validations anchor
the extremes. A Markov changepoint task whose oracle defeats unordered
register-wise merges is solved by a tree-plus-FNO instance of the
neural-operator realization layer
(Appendices~\ref{sec:markov-interlude}--\ref{sec:markov-walkthrough}),
and running the framework with a deterministic register-wise merge
recovers HyperLogLog byte-for-byte, with a fully end-to-end learned
version tracking the classical sketch within a small multiple of its
theoretical relative standard error (Appendix~\ref{sec:hll-parity}).
Appendix~\ref{app:operator-overlap} records the
formal overlap: transformer stacks sit inside equation-(6) neural
operators, mergeable sketches sit inside exact local-law operators, and
the certified intersection is the mergeable neural-operator subspace.
The manifesto application then takes the same machinery to a real
measurement target at two supervision levels. With document labels
only, a prompt-only C-Tree over one open-weight model matches a
three-frontier-LLM ensemble on the Benoit benchmark and stays
chunk-invariant under iterative prompt optimization
(Sections~\ref{sec:benoit-headline}--\ref{sec:rile-fg-ladder}); the
empirical claim is local preservation, not blind faith in a
full-context read. With the corpus's quasi-sentence expert codes
supplying a gold oracle at every node, a trained neural-operator merge
composes on every signal-dense policy dimension, and the failures land
exactly where the label system predicts them: sparse dimensions,
constructs the codes do not determine, objectives that do not price
the local laws
(Sections~\ref{sec:manifesto-qsentence}--\ref{sec:manifesto-lessons}).
The main finding: for
tasks whose structure aligns with the local laws, reliable measurement
on long documents becomes a question of local architecture and audit
budget rather than of full-document annotation alone. The trade
restricts the framework's scope to locally-decomposable oracles, and it
pays off in the many social-scientific settings where such
decomposition is natural.

The Semantic Forests companion extends these ideas to orchestration
and deployment at system scale.
